% !TeX root = ../../main.tex
\section{动态精度控制的数值与执行约束}
\subsection{真实残差与递推残差的关系}
有限精度迭代中，递推残差由前一轮残差减去近似乘积得到，而真实残差由当前近似解重新计算。定义残差间隙 $g_k=(b-Ax_k)-r_k$，忽略其他舍入项，仅考虑 SpMV 误差 $e_k=\widetilde q_k-Ap_k$，可得
\begin{equation}
 g_{k+1}=g_k+\alpha_ke_k.
\end{equation}
其来源是将 $x_{k+1}=x_k+\alpha_kp_k$和 $r_{k+1}=r_k-\alpha_k(Ap_k+e_k)$代入定义。由三角不等式得到
\begin{equation}
 \norm{g_N}\leq\norm{g_0}+\sum_{k=0}^{N-1}|\alpha_k|\norm{e_k}.
\end{equation}
这是忽略其他浮点误差的解释性界，不是完整浮点 CG 的误差分析；实际步长和向量更新也会引入舍入。

该关系说明，即使每轮误差都较小，也需要关注其沿迭代累积的影响。递推残差趋近零时，真实残差可能被间隙限制。结束后以 FP64 重算真实残差，可独立检验递推停止条件所对应的数值质量。若将这种检查加入在线系统，必须把额外 SpMV 成本计入部署时间。

\subsection{状态归一化的收益与信息损失}
归一化输入使网络更关注方向和局部分布，而不受搜索向量整体尺度变化支配。这样可以减轻训练中不同迭代数值范围的差异，并为跨尺度迁移提供可能。但不同绝对幅值的片段归一化后可能相同，其实际误差敏感性却不一定相同。

因此，归一化是一种特征选择，而非无损变换。可将片段范数、残差阶段或历史动作加入状态作为扩展，但这会增加推理输入和训练复杂度。比较不同状态设计时，应固定执行基础、网络预算和训练协议，避免把更大网络的优势误归因于特定数值特征。

\subsection{动作粒度与控制复杂度}
列块宽度 $C$决定策略调用数量，约为 $\lceil n/C\rceil$。较小 $C$提供更细的局部精度选择，却增加状态提取、动作存储和推理成本；较大 $C$降低控制开销，但可能让局部敏感区域迫使整个块使用高精度。Tile 行高 $R$又影响物理执行与元数据，两个参数的作用不同。

\begin{table}[htbp]\centering\caption{EvoCG 空间与时间粒度的作用}
\begin{tabularx}{\textwidth}{p{2.5cm}XX}\toprule
粒度参数&主要影响&需要联动观察的指标\\\midrule
列块宽度 $C$&动作覆盖范围和状态维数&精度比例、推理时间、收敛\\
行块高度 $R$&任务划分和 Tile 数量&负载分布、索引开销、补齐\\
每行块线程束数&并发 Tile 执行&占用率、同步与资源压力\\
复用间隔 $K^*$&策略进入关键路径的频率&升精度延迟、迭代数、控制成本\\\bottomrule
\end{tabularx}\end{table}

本实现固定64 列，与线程束协同遍历和网络结构相匹配。若改变 $C$，不仅需要重训练网络，还可能改变列指针长度、向量化访问和内核映射。因此，参数搜索应围绕完整系统进行，不能单独选择最好的网络输入长度后假设硬件成本不变。

\subsection{统一索引下的容量成本}
设逻辑非零数为 $z$，共享行索引每项 $b_i$字节。FP64、FP32、BF16 三份数值合计14 字节，单份 FP64 为8 字节。忽略 Tile 描述符与补齐，数值加索引的存储比近似为
\begin{equation}
 \rho_{\mathrm{base}}=\frac{14+b_i}{8+b_i}.
\end{equation}
例如以4 字节索引作说明，比例为1.5，而不是数值字段的1.75。实际 Tile-CSC 还包含列指针、Tile 行指针、描述符和填充，因此测得比例可能明显更高。

若补齐使物理非零槽位数为 $z_p\geq z$，这些数值流和索引成本应按 $z_p$计算。对于大量只含一个非零的列，补齐到2 会造成较高相对开销；若非零元素分散到更多 Tile，每 Tile 固定元数据又会放大总占用。实验离群矩阵与这一机制相符，但精确归因仍需逐项统计字节数。

\subsection{策略复用的保守性与失效条件}
精度动作在相邻迭代保持不变，才使区间复用具有收益。如果求解状态突然进入数值敏感阶段，长复用窗口可能错过及时升精度。原方案通过离线轨迹控制延迟升精度的比例，这是经验约束，而非对任意在线轨迹的防护。

可用于分析的指标包括每个列块的动作切换次数、升降精度比例、连续保持长度以及延迟切换对应的迭代增量。只报告平均推理占比不能说明复用策略是否合理。特别是极短求解任务，一次错误复用就可能占据很大比例，需要与长迭代问题分别分析。

部署时若右端项或初值改变，即使矩阵未变，状态轨迹也可能改变。此时是否继续使用相同 $K^*$需要验证。保守回退可缩短间隔或使用 FP64，但相应成本应进入系统评价，不能将失败检测与回退视为额外免费步骤。

\subsection{训练随机性与策略选择偏差}
强化学习训练包含动作采样和优化随机性，单次训练得到的策略不一定代表稳定效果。完整复现应记录随机种子、回合预算、检查点选择准则以及用于选择策略的验证负载。若仅报告若干训练中的最佳在线时间，可能引入选择偏差。

对矩阵专属策略，应说明训练和评价是否使用相同右端项；对跨矩阵策略，应保持测试矩阵未参与调参。本文所整理的原结果具有明确矩阵集合和默认训练设置，多随机种子下的稳定性和置信区间仍有待补充评估。后续重复实验应在固定协议下报告成功率与时间分布，而非只报告最优一次。

\subsection{失败样本的计分方法}
当某方法未在迭代上限内完成求解，其运行时间是失败耗时而非成功求解时间。把失败样本从全部平均中直接删除，会使保留子集偏向容易问题；把上限时间当作完成时间，又会混淆质量要求。因此，应先报告成功数量，再在明确的共同成功集合上比较时间。

\begin{table}[htbp]\centering\caption{求解实验中不同结果状态的报告方式}
\begin{tabularx}{\textwidth}{p{2.8cm}XX}\toprule
结果状态&质量信息&性能信息\\\midrule
按条件成功&递推残差及真实残差&成功求解时间、迭代次数\\
达到迭代上限&注明未满足停止条件&仅报告失败耗时，不计为成功时间\\
出现非有限数&记录首次发生位置&保留故障前耗时及状态\\
实现不支持&说明矩阵或格式约束&从该实现比较集合中明确排除\\
策略回退成功&分别记录回退前后质量&计入检测、回退与重算总时间\\\bottomrule
\end{tabularx}\end{table}

原论文对 Mille 系列的比较采用其发布实现成功完成的矩阵集合，跨矩阵迁移又包含9 个失败案例。二者都需要保留集合信息。本文将相同执行基础下的矩阵专属结果作为精度策略的主要证据，将迁移结果作为有边界的补充证据。

\subsection{预条件迁移与求解器结构变化}
Jacobi 预条件改变搜索方向的数值分布，但保持较简单的逐元素应用，因此提供一种受控的策略迁移测试。更复杂预条件器可能引入额外稀疏三角求解、不同精度误差和更高单轮成本，这时 SpMV 在总时间中的占比会改变。

若预条件使迭代次数显著减少，策略训练与推理成本的相对比例可能增加。若预条件计算成为主要开销，仅加速 SpMV 的收益又受 Amdahl 效应限制。因此，后续扩展需同时测量精度控制的数值效果和算子耗时构成，不能从 Jacobi 实验直接推断更复杂预条件器上的同等收益。
